To establish that the proposed approach is feasible and worth pursuing, we have
implemented a proof-of-concept of McpRisk and run an initial evaluation. This
chapter reports those \emph{preliminary} results; a fuller evaluation is part of
the planned work in \Cref{chapter:discussion_and_conclusions}. Because risk
banding is intrinsically subjective, we do not evaluate against a single
``correct'' answer but triangulate across four kinds of ground truth, and we
report both the results that favor the approach and those that do not.

\section{Setup}
\label{sec:experiments:setup}
\addcontentsline{tocheb}{section}{\protect\numberline{\secnumforhebrewtoc}{הגדרות הניסוי}}

The evaluation uses four datasets. First, \textbf{seven realistic MCP
deployments} --- corporate, legal, medical, and media filesystems, and two SQLite
databases --- each seeded with crown-jewel assets (private keys, payslips,
patient records, API tokens); their judge-reviewed static tables ($61$ tools,
over $40$ assets, $5$ \textit{critical} cells) form the design-time ground truth.
Second, an \textbf{independent ten-rater oracle panel} per cell, combining a
hand-authored human heatmap with nine framework baselines (DREAD, CVSS~v3, two
NIST methods, OWASP Risk Rating, OWASP AIVSS, MAESTRO, and two ChatGPT raters);
the scanner's own model is excluded to avoid circularity. Third, a
\textbf{captured-call corpus} of $969$ calls from seeded simulations across
Calendar, GitHub, Slack, a fintech filesystem, and a DevOps database, mixing
benign operators with attackers. Fourth, $150$ labeled calls from three
\textbf{external attack benchmarks} (MCPSecBench, MCP-SafetyBench, MSB) that share
our client~$\to$~server threat model.

All scores were produced by a single local model (Qwen2.5:32B via Ollama,
temperature $0$, seed $0$). Since the output is an ordinal four-band label, we
report \emph{exact} and \emph{within-one-band} agreement, quadratic-weighted
Cohen's $\kappa$, and --- for the detection experiments --- the area under the ROC
curve (AUC). Panel consensus is the upper-median across raters, robust to the
maximalism of frameworks such as DREAD that rate almost everything dangerous.

\section{Static Scoring vs.\ Human Judgment}
\label{sec:experiments:static}
\addcontentsline{tocheb}{section}{\protect\numberline{\secnumforhebrewtoc}{ניקוד סטטי מול שיפוט אנושי}}

The two \emph{primitive} judgments are highly consistent with the reviewed
tables: tool-impact agreement is $59/61$ ($97\%$) and asset-sensitivity
agreement $5/7$ ($71\%$). Agreement on the \emph{composite} band is lower ($34\%$
exact), the first sign that the composite score --- not the primitives --- is the
weak point. Against the independent panel, the crucial context is the
\emph{legitimate-disagreement ceiling}: the human and framework raters agree with
\emph{each other} only $50\%$ (filesystem) and $42\%$ (SQLite) of the time
exactly (\Cref{tab:human}). Risk banding is inherently subjective, and no scorer
can exceed this ceiling. McpRisk reaches $35\%$ exact and $69\%$ within-one
agreement with the consensus --- inside the band of human disagreement, though at
its lower end --- and tends to score \emph{below} the conservative consensus,
i.e.\ it under-alarms rather than over-alarms.

\begin{table}[H]
    \centering
    \caption{Scanner vs.\ the human/framework panel, against the inter-rater
    ceiling. ``Exact'' and ``$\le 1$'' are band-agreement rates.}
    \label{tab:human}
    \begin{tabular}{l c c}
        \hline
        \textbf{Comparison} & \textbf{Exact} & \textbf{$\le 1$ band} \\
        \hline
        Inter-rater ceiling (filesystem) & 50\% & 78\% \\
        Inter-rater ceiling (SQLite)     & 42\% & 75\% \\
        Scanner vs.\ consensus (overall) & 35\% & 69\% \\
        \hline
    \end{tabular}
\end{table}

Scoring the same cells with every baseline and comparing against the human oracle
by quadratic-weighted $\kappa$ (\Cref{tab:headtohead}), McpRisk attains
$\kappa=+0.66$ with $97\%$ within-one agreement --- second only to OWASP AIVSS
and well ahead of CVSS~v3. The maximalist frameworks reach apparent safety
trivially by rating almost everything dangerous (DREAD $\kappa=+0.10$ at a
$100\%$ over-block rate); McpRisk's $56\%$ over-block matches AIVSS while
preserving discrimination, the intended behavior of the risk pyramid.

\begin{table}[H]
    \centering
    \caption{Agreement with the human oracle (pooled). QW-$\kappa$ is
    quadratic-weighted Cohen's $\kappa$; over-block is the fraction rated
    high/critical.}
    \label{tab:headtohead}
    \begin{tabular}{l c c c}
        \hline
        \textbf{Scorer} & \textbf{QW-$\kappa$} & \textbf{Within-1} & \textbf{Over-block} \\
        \hline
        OWASP AIVSS             & $+0.69$ & --    & 56\% \\
        \textbf{McpRisk (ours)} & $+0.66$ & 97\%  & 56\% \\
        CVSS v3                 & $+0.36$ & --    & --   \\
        ChatGPT (plain)         & $+0.23$ & --    & 98\% \\
        DREAD                   & $+0.10$ & --    & 100\% \\
        \hline
    \end{tabular}
\end{table}

\section{Robustness, Request-Time, and Out-of-Distribution Detection}
\label{sec:experiments:dynamic}
\addcontentsline{tocheb}{section}{\protect\numberline{\secnumforhebrewtoc}{יציבות, זמן ריצה וזיהוי מחוץ להתפלגות}}

The approach's most significant weakness is the \emph{stability} of the composite
score. Perturbing each of the $1{,}700$ scored cells' inputs by $\pm1$, one at a
time, changes the band for $85\%$ of cells, with the fragility clustered at the
\textit{medium} threshold. The multiplicative formula with fixed integer
cut-points is thus unstable precisely where a gate's behavior hinges --- the
primary target of the planned work.

On the captured-call corpus, McpRisk produces a scorable band for $87\%$ of a
labeled subset, abstaining honestly on the rest, and reaches an AUC of $0.64$
against the sessions' authored-intent labels. These labels are a \emph{weak}
ground truth, so the numbers are diagnostic; the informative failure mode is that
$27$ malicious \emph{read-only} exfiltration calls score only low/medium, since
the inherent cell under-weights pure exfiltration --- exactly the gap the planned
contextual scorer must close. The out-of-distribution result is the most
encouraging (\Cref{tab:external}): on the external benchmarks, content-aware
scoring separates attacks (LLM-judge AUC $0.68$) while the argument-blind
capability frameworks CVSS and AIVSS score $0.52$ --- essentially no skill,
because they cannot see the arguments that make a call malicious. This is direct
evidence for the proposal's premise that risk must be assessed from the request
content, not the tool's capability alone.

\begin{table}[H]
    \centering
    \caption{Attack-vs-valid detection on $150$ external-benchmark calls.}
    \label{tab:external}
    \begin{tabular}{l c c}
        \hline
        \textbf{Scorer} & \textbf{AUC} & \textbf{Recall @ 10\% FPR} \\
        \hline
        LLM judge (content-aware) & 0.68 & 42\% \\
        Keyword (content-aware)   & 0.67 & 38\% \\
        CVSS (capability-only)    & 0.52 & 0\%  \\
        AIVSS (capability-only)   & 0.52 & 0\%  \\
        Random baseline           & 0.48 & 11\% \\
        \hline
    \end{tabular}
\end{table}

Finally, an independent review subsystem audits the pipeline: a deterministic
verifier passes all $8/8$ correctness checks and a methodology judge confirms all
$7/7$ integrity claims (no data leakage, LLM-only with no fabrication, honest
handling of unresolved calls), while a results judge returns a sober verdict of
``adequate, with concern'' --- the pipeline is correct and honest, but its
agreement with the human consensus is only modestly above chance. These
preliminary findings, and the weaknesses they expose, motivate the research plan
of the next chapter.
